\begin{resumo}
Este projeto tem como objetivo desenvolver um micromouse aut\^onomo de baixo custo, capaz de navegar em um labirinto desconhecido sem interven\c{c}\~ao humana, atendendo \`as regras da competi\c{c}\~ao IEEE. A solu\c{c}\~ao integra hardware embarcado, sensoriamento e controle de motores, com algoritmos de navega\c{c}\~ao e mapeamento em tempo real, al\'em de um sistema web para telemetria. O rob\^o sera concebido com componentes acess\'iveis no mercado nacional, priorizando confiabilidade e facilidade de replica\c{c}\~ao. Entre os requisitos principais est\~ao o movimento aut\^onomo, detec\c{c}\~ao de paredes, evita\c{c}\~ao de colis\~oes e envio de dados de telemetria com baixa lat\^encia. A arquitetura de software prev\^e comunica\c{c}\~ao sem fio, armazenamento das trajet\'orias e visualiza\c{c}\~ao de m\'etricas de desempenho. O resultado esperado \`e uma plataforma did\'atica e funcional de algoritmos em ambiente controlado.

\vspace{\onelineskip}
\noindent
\textbf{Palavras-chaves}: \imprimirpalavrachaveum, \imprimirpalavrachavedois
\end{resumo}

% \begin{resumo}
%  O resumo deve ressaltar o objetivo, o método, os resultados e as conclusões 
%  do documento. A ordem e a extensão
%  destes itens dependem do tipo de resumo (informativo ou indicativo) e do
%  tratamento que cada item recebe no documento original. O resumo deve ser
%  precedido da referência do documento, com exceção do resumo inserido no
%  próprio documento. (\ldots) As palavras-chave devem figurar logo abaixo do
%  resumo, antecedidas da expressão Palavras-chave:, separadas entre si por
%  ponto e finalizadas também por ponto. O texto pode conter no mínimo 150 e 
%  no máximo 500 palavras, é aconselhável que sejam utilizadas 200 palavras. 
%  E não se separa o texto do resumo em parágrafos.

%  \vspace{\onelineskip}
    
%  \noindent
%  \textbf{Palavras-chave}: latex. abntex. editoração de texto.
% \end{resumo}
